\chapter{剪切}

\section{受力与变形特点}

\subsection{受力特点}
构件受一对大小相等、方向相反、作用线相距很近且垂直于杆轴线的横向外力作用。
\begin{itemize}
    \item \textbf{受力：} 外力平行、方向相反，但不在一条直线上。
    \item \textbf{变形特点：} 沿剪切面剪断（相对错动变形）。
    \item \textbf{破坏形式：} 剪切破坏、挤压破坏（在应力传递处导致）。
\end{itemize}

\section{剪切的实用计算}

\subsection{内力分析}
通过截面法，取上半部分分析：
\begin{equation}
    \sum X_i = 0 \implies F - F_s = 0 \implies F_s = F
\end{equation}
其中 $F_s$ 为剪切面上的内力，称为\textbf{剪力}。

\subsection{切应力计算}
假设切应力在剪切面上均匀分布，则切应力计算公式为：
\begin{equation}
    \tau = \frac{F_s}{A}
\end{equation}
其中 $A$ 为剪切面积。

\subsection{切应力强度条件}
为了保证构件不发生剪切破坏，需满足：
\begin{equation}
    \tau = \frac{F_s}{A} \le [\tau]
\end{equation}
其中 $[\tau]$ 为材料的许用切应力。

\section{挤压的实用计算}

\subsection{挤压应力公式}
挤压发生在连接件与被连接件的接触表面。其实用计算公式为：
\begin{equation}
    \sigma_{bs} = \frac{F_{bs}}{A_{bs}}
\end{equation}
\begin{itemize}
    \item $F_{bs}$：挤压力。
    \item $A_{bs}$：有效挤压面积。
\end{itemize}

\subsection{挤压面积的确定}
对于圆柱形接触面（如螺栓、销钉），$A_{bs}$ 采用\textbf{正投影面积}。
\begin{itemize}
    \item \textbf{思考：} 为什么不采用实际接触面积（半圆弧面）？
    \item \textbf{结论：} 采用投影面积计算出的 $\sigma_{bs}$ 更大，结果更趋于安全，计算也更简便。
\end{itemize}

